% GENERATED FILE. Edit canonical lessons, then run npm run generate:books.
% canonical-source-hash: fnv1a64:8960535161010cc3
% canonical-lessons: PA-C225-pasarna, PA-C225-jagauna, PA-C225-vihla, PA-C225-letna, PA-C225-kahla

\chapter{Waking Up Early}
\label{ch:pa-a2-waking-up-early}
\hlchaptermodality{225}

\begin{chapteropening}
\textbf{By the end of this chapter:} \emph{I can say when I wake and when I am free.}

Complete the last lesson of chapter 225: I can say when I wake and when I am free.
\end{chapteropening}

\section[\texorpdfstring{\pa{ਪਸਾਰਨਾ} (pasārnā)}{ਪਸਾਰਨਾ (pasārnā)}]{\pa{ਪਸਾਰਨਾ} (pasārnā) — to stretch out (arms or legs)}

\label{lesson:PA-C225-pasarna}



\begin{quote}
Before the new one: say the Punjabi for to take out, then the Punjabi for perhaps.
\end{quote}

\subsection*{You'll want to know: \pa{ਪਸਾਰਨਾ}}
\textbf{\pa{ਪਸਾਰਨਾ}} — \emph{pasārnā} — \textquotedblleft{}to stretch out (arms or legs)\textquotedblright{}.

\begin{quote}
\pa{ਸਵੇਰੇ} \pa{ਮੈਂ} \pa{ਮੰਜੇ} \pa{ਉੱਤੇ} \pa{ਲੱਤਾਂ} \pa{ਪਸਾਰਦਾ} \pa{ਹਾਂ।}
\end{quote}

\begin{quote}
\emph{savere maiṁ manje utte lattāṁ pasārdā hāṁ.}
\end{quote}

\begin{quote}
In the morning I stretch out my legs on the cot.
\end{quote}

Past (I …): \pa{ਮੈਂ} \pa{ਲੱਤਾਂ} \pa{ਪਸਾਰੀਆਂ}, \emph{maiṁ lattāṁ pasārīāṁ}.

\begin{cousinweb}
\textbf{\pa{ਪਸਾਰਨਾ}} (\emph{pasārnā}) is inherited from Sanskrit \emph{prasārayati}, \textquotedblleft{}stretches out, extends\textquotedblright{}; CDIAL lists Punjabi \emph{pasārṇā}, \textquotedblleft{}to stretch out\textquotedblright{} (CDIAL 8838). CDIAL adds that the Punjabi form may have been taken from the central languages.
\end{cousinweb}

\begin{grammarlens}[title={What you've built}]
One more word for saying what you do.
\end{grammarlens}

\subsection*{Your turn}
\begin{itemize}
\raggedright
  \item \emph{Say it:} \emph{pasārnā}
  \item \emph{Say it:} \emph{pasārnā}, once more
  \item \emph{Say it:} \emph{savere maiṁ manje utte lattāṁ pasārdā hāṁ.} — In the morning I stretch out my legs on the cot.
  \item \emph{Recall:} say the Punjabi for to take out, then the Punjabi for perhaps, then say \emph{pasārnā} again
\end{itemize}

\begin{tcolorbox}[breakable,colback=teal!4,colframe=teal!35!black,title={Before you move on}]
What does \pa{ਪਸਾਰਨਾ} mean? (\textquotedblleft{}To stretch out (arms or legs)\textquotedblright{}.) Say it once more.
\end{tcolorbox}
\section[\texorpdfstring{\pa{ਜਗਾਉਣਾ} (jagāuṇā)}{ਜਗਾਉਣਾ (jagāuṇā)}]{\pa{ਜਗਾਉਣਾ} (jagāuṇā) — to wake (someone) up}

\label{lesson:PA-C225-jagauna}



\begin{quote}
Before the new one: say the Punjabi for perhaps, then the Punjabi for to stretch out.
\end{quote}

\subsection*{You'll want to know: \pa{ਜਗਾਉਣਾ}}
\textbf{\pa{ਜਗਾਉਣਾ}} — \emph{jagāuṇā} — \textquotedblleft{}to wake (someone) up\textquotedblright{}.

\begin{quote}
\pa{ਮੈਂ} \pa{ਸਵੇਰੇ} \pa{ਛੇ} \pa{ਵਜੇ} \pa{ਬੱਚਿਆਂ} \pa{ਨੂੰ} \pa{ਜਗਾਉਂਦੀ} \pa{ਹਾਂ।}
\end{quote}

\begin{quote}
\emph{maiṁ savere chhe vaje bachchiāṁ nū̃ jagāundī hāṁ.}
\end{quote}

\begin{quote}
I wake the children at six in the morning.
\end{quote}

Past (I …): \pa{ਮੈਂ} \pa{ਅਮਨ} \pa{ਨੂੰ} \pa{ਜਗਾਇਆ}, \emph{maiṁ Aman nū̃ jagāiā}.

\begin{cousinweb}
\textbf{\pa{ਜਗਾਉਣਾ}} (\emph{jagāuṇā}) is the causative partner of \pa{ਜਾਗਣਾ}, built on a verb inherited from Sanskrit \emph{jāgrati}, \textquotedblleft{}are awake\textquotedblright{}; CDIAL lists Punjabi \emph{jagāuṇā} under it (CDIAL 5175).
\end{cousinweb}

\begin{grammarlens}[title={What you've built}]
One more word for saying what you do.
\end{grammarlens}

\subsection*{Your turn}
\begin{itemize}
\raggedright
  \item \emph{Say it:} \emph{jagāuṇā}
  \item \emph{Say it:} \emph{jagāuṇā}, once more
  \item \emph{Say it:} \emph{maiṁ savere chhe vaje bachchiāṁ nū̃ jagāundī hāṁ.} — I wake the children at six in the morning.
  \item \emph{Recall:} say the Punjabi for perhaps, then the Punjabi for to stretch out, then say \emph{jagāuṇā} again
  \item \emph{Write it:} copy \textbf{\pa{ਜਗਾਉਣਾ}} once, checking the whole word
\end{itemize}

\begin{tcolorbox}[breakable,colback=teal!4,colframe=teal!35!black,title={Before you move on}]
What does \pa{ਜਗਾਉਣਾ} mean? (\textquotedblleft{}To wake (someone) up\textquotedblright{}.) Say it once more.
\end{tcolorbox}
\section[\texorpdfstring{\pa{ਵਿਹਲਾ} (vihlā)}{ਵਿਹਲਾ (vihlā)}]{\pa{ਵਿਹਲਾ} (vihlā) — free; not busy}

\label{lesson:PA-C225-vihla}



\begin{quote}
Before the new one: say the Punjabi for to stretch out, then the Punjabi for to wake (someone) up.
\end{quote}

\subsection*{You'll want to know: \pa{ਵਿਹਲਾ}}
\textbf{\pa{ਵਿਹਲਾ}} — \emph{vihlā} — \textquotedblleft{}free; not busy\textquotedblright{}.

\begin{quote}
\pa{ਸ਼ਾਮ} \pa{ਨੂੰ} \pa{ਮੈਂ} \pa{ਵਿਹਲਾ} \pa{ਹੁੰਦਾ} \pa{ਹਾਂ।}
\end{quote}

\begin{quote}
\emph{shām nū̃ maiṁ vihlā hundā hāṁ.}
\end{quote}

\begin{quote}
In the evening I am free.
\end{quote}

\begin{cousinweb}
\textbf{\pa{ਵਿਹਲਾ}} (\emph{vihlā}) is inherited from a reconstructed \emph{*vēlya}, \textquotedblleft{}at leisure\textquotedblright{}; CDIAL lists Punjabi \emph{vihlā}, \textquotedblleft{}vacant, at leisure\textquotedblright{} (CDIAL 12118).
\end{cousinweb}

\begin{grammarlens}[title={What you've built}]
One more word for saying what you do.
\end{grammarlens}

\subsection*{Your turn}
\begin{itemize}
\raggedright
  \item \emph{Say it:} \emph{vihlā}
  \item \emph{Say it:} \emph{vihlā}, once more
  \item \emph{Say it:} \emph{shām nū̃ maiṁ vihlā hundā hāṁ.} — In the evening I am free.
  \item \emph{Recall:} say the Punjabi for to stretch out, then the Punjabi for to wake (someone) up, then say \emph{vihlā} again
  \item \emph{Write it:} copy \textbf{\pa{ਵਿਹਲਾ}} once, checking the whole word
\end{itemize}

\begin{tcolorbox}[breakable,colback=teal!4,colframe=teal!35!black,title={Before you move on}]
What does \pa{ਵਿਹਲਾ} mean? (\textquotedblleft{}Free; not busy\textquotedblright{}.) Say it once more.
\end{tcolorbox}
\section[\texorpdfstring{\pa{ਲੇਟਣਾ} (leṭṇā)}{ਲੇਟਣਾ (leṭṇā)}]{\pa{ਲੇਟਣਾ} (leṭṇā) — to lie down}

\label{lesson:PA-C225-letna}



\begin{quote}
Before the new one: say the Punjabi for to wake (someone) up, then the Punjabi for not busy.
\end{quote}

\subsection*{You'll want to know: \pa{ਲੇਟਣਾ}}
\textbf{\pa{ਲੇਟਣਾ}} — \emph{leṭṇā} — \textquotedblleft{}to lie down\textquotedblright{}.

\begin{quote}
\pa{ਦੁਪਹਿਰ} \pa{ਨੂੰ} \pa{ਮੈਂ} \pa{ਮੰਜੇ} \pa{ਉੱਤੇ} \pa{ਲੇਟਦਾ} \pa{ਹਾਂ।}
\end{quote}

\begin{quote}
\emph{dupahir nū̃ maiṁ manje utte leṭdā hāṁ.}
\end{quote}

\begin{quote}
In the afternoon I lie down on the cot.
\end{quote}

Past (I …): \pa{ਮੈਂ} \pa{ਮੰਜੇ} \pa{ਉੱਤੇ} \pa{ਲੇਟਿਆ}, \emph{maiṁ manje utte leṭiā}.

\begin{cousinweb}
\textbf{\pa{ਲੇਟਣਾ}} (\emph{leṭṇā}) is inherited from Sanskrit \emph{lēṭyati}, \textquotedblleft{}sleeps\textquotedblright{}; CDIAL lists Punjabi \emph{leṭṇā}, \textquotedblleft{}to lie down, roll\textquotedblright{} (CDIAL 11109).
\end{cousinweb}

\begin{grammarlens}[title={What you've built}]
One more word for saying what you do.
\end{grammarlens}

\subsection*{Your turn}
\begin{itemize}
\raggedright
  \item \emph{Say it:} \emph{leṭṇā}
  \item \emph{Say it:} \emph{leṭṇā}, once more
  \item \emph{Say it:} \emph{dupahir nū̃ maiṁ manje utte leṭdā hāṁ.} — In the afternoon I lie down on the cot.
  \item \emph{Recall:} say the Punjabi for to wake (someone) up, then the Punjabi for not busy, then say \emph{leṭṇā} again
  \item \emph{Write it:} copy \textbf{\pa{ਲੇਟਣਾ}} once, checking the whole word
\end{itemize}

\begin{tcolorbox}[breakable,colback=teal!4,colframe=teal!35!black,title={Before you move on}]
What does \pa{ਲੇਟਣਾ} mean? (\textquotedblleft{}To lie down\textquotedblright{}.) Say it once more.
\end{tcolorbox}
\section[\texorpdfstring{\pa{ਕਾਹਲਾ} (kāhlā)}{ਕਾਹਲਾ (kāhlā)}]{\pa{ਕਾਹਲਾ} (kāhlā) — in a hurry; hasty}

\label{lesson:PA-C225-kahla}



\begin{quote}
Before the new one: say the Punjabi for not busy, then the Punjabi for to lie down.
\end{quote}

\subsection*{You'll want to know: \pa{ਕਾਹਲਾ}}
\textbf{\pa{ਕਾਹਲਾ}} — \emph{kāhlā} — \textquotedblleft{}in a hurry; hasty\textquotedblright{}.

\begin{quote}
\pa{ਸਵੇਰੇ} \pa{ਮੈਂ} \pa{ਬਹੁਤ} \pa{ਕਾਹਲਾ} \pa{ਹੁੰਦਾ} \pa{ਹਾਂ।}
\end{quote}

\begin{quote}
\emph{savere maiṁ bahut kāhlā hundā hāṁ.}
\end{quote}

\begin{quote}
In the mornings I am in a great hurry.
\end{quote}

\begin{cousinweb}
\textbf{\pa{ਕਾਹਲਾ}} (\emph{kāhlā}) is inherited: CDIAL sets it under Sanskrit \emph{kāhala}, \textquotedblleft{}speaking indistinctly, faint-hearted\textquotedblright{}, and lists Punjabi \emph{kāhlā}, \textquotedblleft{}impatient, hasty\textquotedblright{} (CDIAL 3142); CDIAL does not give the everyday sense \textquotedblleft{}in a hurry\textquotedblright{} as such.
\end{cousinweb}

\begin{grammarlens}[title={What you've built}]
One more word for saying what you do.
\end{grammarlens}

\subsection*{Your turn}
\begin{itemize}
\raggedright
  \item \emph{Say it:} \emph{kāhlā}
  \item \emph{Say it:} \emph{kāhlā}, once more
  \item \emph{Say it:} \emph{savere maiṁ bahut kāhlā hundā hāṁ.} — In the mornings I am in a great hurry.
  \item \emph{Recall:} say the Punjabi for not busy, then the Punjabi for to lie down, then say \emph{kāhlā} again
  \item \emph{Write it:} copy \textbf{\pa{ਕਾਹਲਾ}} once, checking the whole word
\end{itemize}

\begin{tcolorbox}[breakable,colback=teal!4,colframe=teal!35!black,title={Before you move on}]
What does \pa{ਕਾਹਲਾ} mean? (\textquotedblleft{}In a hurry; hasty\textquotedblright{}.) Say it once more.
\end{tcolorbox}
